\documentclass[a4paper,10pt]{article}
\newcommand\subdir{/home/koen/LaTeX-setup/}
\input{\subdir/LaTeX-files/imports.tex}
\input{\subdir/LaTeX-files/master-internship.tex}
\usepackage{\subdir/LaTeX-files/aastex}
\usepackage{natbib}
\bibliographystyle{\subdir/LaTeX-files/astroadstitle.bst}

%Set the title, Author and date
\title{Notes Week 35}
\author{Koen Essers}
\tutorial{Master Internship}
\date{\today}

%Set the header style
\header{\tutorialstring}

%Begin the document
\begin{document}
\setuplayout
\section{Different Masses}

I want to see if there is some predictable structure in the plot below:

\begin{figure}[ht]
	\marginfignote{0}{The dots represent the mass of the TPAGB star at the start of a certain thermal pulse.
		The exact thermal pulse numbers are irrelevant here.
		It is only important to look at the lowest two dot.
		These dot represent the mass of the envelope during the \emph{last} two thermal pulse.
		Clearly, if we interact during
		the last \emph{interpulse} phase, RLOF will occur when the star has an envelope mass between these two values.
	}
	\marginfignote{17}{The figure above shows that for the models that interact during the interpulse, the envelope mass is approximately \(0.5\; M_\odot \).
		This figure shows that this is totally in agreement with the values that we expect are possible.
		The envelope mass of the \(1.4\; M_\odot \) star during the last thermal pulse is not very low compared to the other stars.
	}
	\centering
	\incplot{w33-mass-at-TP}
\end{figure}

\begin{figure}[ht]
	\marginfignote{0}{During the last thermal pulses, there is quite a strong correlation between the
		mass and the interpulse duration.
		However (unfortunately), during the earlier thermal pulses, this relation is a bit less clear.
	}
	\centering
	\incplot{w35-interpulse-time}
\end{figure}

Unfortunately, the changing interpulse duration \emph{already} makes it difficult to predict the
mass of the envelope during the thermal pulses.
This is even without taking into account that
the wind mass loss is also not constant throughout the evolution.

I am interested to see how drastically the envelope mass changes between masses.
The first figure makes it seem like it is very chaotic.
I want to simulate some more single stars
between some masses to see if the chaotic nature is really part of the TP-AGB evolution, \emph{or}, if it is
just an effect of the rough sampling or \(0.1\; M_\odot \) I am using.

I choose to simulate the masses \(1.9\) to \(2.1\) with a delta of \(0.01\; M_\odot \).
In the first figure,
we see that these masses have roughly the same envelope mass during the final thermal pulse.
However, they all have a different amount of thermal pulses, so the evolution \emph{must} be quite different.

% \begin{figure}[ht]
% 	\centering
% 	\incplot{w35-roche-radii}
% \end{figure}
%
% \nextpage
\begin{figure}[ht]
	\centering
	\incplot{w35-roche-radii-2}
\end{figure}

\begin{figure}[ht]
	\marginfignote{3}{This plot shows the final period of a grid of models
		with \(q=0.85,\;\varepsilon=0.1\) and \(f_\text{delta}=0.1\).}
	\marginfignote{8}{Unfortunately, the resolution of the models in the
		\(M_\text{TPAGB}\)--\(R_\text{RL}\) plane is not great.
		I color
		the plane using cubic interpolation between the datapoints.
		The contours are definitely not precise, and not too much importance should be attributed to the contours.
	}
	\marginfignote{18}{These models have the \emph{lowest} \(f_{\delta }\) values.
		This means these models have the \emph{largest} periods.
	}
	\centering
	\incplot{w35-mass-1}
\end{figure}

\begin{figure}[ht]
	\marginfignote{3}{This plot shows the final period of a grid of models
		with \(q=0.85,\;\varepsilon=0.1\) and \(f_\text{delta}=0.2\).}
	\marginfignote{8}{As expected, the higher value of \(f_{\delta }=0.2\) produces lower final periods.
		However, there is also a slightly different effect that can be seen when comparing this plot with the one above:
		The mass becomes less important, while the initial Roche lobe radius stays relevant.
	}
	\centering
	\incplot{w35-mass-2}
\end{figure}

\begin{figure}[ht]
	\marginfignote{3}{This plot shows the final period of a grid of models
		with \(q=0.85,\;\varepsilon=0.1\) and \(f_\text{delta}=0.3\).}
	\marginfignote{8}{The same trends are visible for the models with \(f_{\delta } = 0.3 \).
		Now, the mass \emph{really} seems like it is
		not relevant anymore.
	}
	\marginfignote{16}{Obviously, things would change if we would take into account the models that
		do not undergo RLOF.}
	\centering
	\incplot{w35-mass-3}
\end{figure}

\nextpage

Let me show how these wind only models would complete the plane on the right.
\begin{figure}[ht]
	\marginfignote{3}{This plot shows a much broader plane of the final period of a grid of models
		with \(q=0.85,\;\varepsilon=0.1\) and \(f_\text{delta}=0.1\).}
	\marginfignote{10}{Clearly the models on the right follow quite a different trend
		then the RLOF models.
		For the wind models, we see that the final period
		follows the same trend as the initial period, where the longest periods correspond
		to the \emph{largest} initial Roche lobe radius and the \emph{least heavy} donor stars.
	}
	\centering
	\incplot{w35-mass-5}
\end{figure}

\begin{figure}[ht]
	\marginfignote{3}{This plot shows a much broader plane of the final period of a grid of models
		with \(q=0.85,\;\varepsilon=0.1\) and \(f_\text{delta}=0.3\).}
	\marginfignote{10}{For the \(f_{\delta } = 0.3\) models, we see that the
		discrepancy between the two methods is much less obvious.
		We see a
		slight shift from pretty much vertical contour lines to slightly diagonal contour lines.
	}
	\marginfignote{19}{Obviously, for the wind only models, there is no difference between
		the \(f_{\delta } =0.3 \) and \(f_{\delta } =0.1 \) models.}
	\centering
	\incplot{w35-mass-4}
\end{figure}

Now, obviously I have only shown \(q=0.85\).
However, the trends still exist for
different mass ratios.
In the appendix, I show \(q=0.7\) as well.

Another exciting thing is that I can probably make a much nicer version of this
plot when my new grid finishes simulating.

Let me show the [s/Fe] abundances for \(q=0.7\).

\begin{figure}[ht]
	\centering
	\incplot{w35-mass-s-1}
\end{figure}

I think this looks great.
Again, we could use the simple binary simulation to
extend the models to higher initial Roche lobe radii, however, this is really not that relevant.
We know these models with very gradually produce lower accreted abundances.
This means
we can easily create barium stars with wind alone, but this is not what this project is about.

\nextpage
\bibliography{master.bib}
\appendix
\section{Additional plots}
\begin{figure}[ht]
	\marginfignote{3}{This plot shows the final period of a grid of models
		with \(q=0.85,\;\varepsilon=0.1\) and \(f_\text{delta}=0.1\).}
	\centering
	\incplot{w35-mass-1-q07}
\end{figure}

\begin{figure}[ht]
	\marginfignote{3}{This plot shows the final period of a grid of models
		with \(q=0.85,\;\varepsilon=0.1\) and \(f_\text{delta}=0.2\).}
	\centering
	\incplot{w35-mass-2-q07}
\end{figure}

\begin{figure}[ht]
	\marginfignote{3}{This plot shows the final period of a grid of models
		with \(q=0.85,\;\varepsilon=0.1\) and \(f_\text{delta}=0.3\).}
	\centering
	\incplot{w35-mass-3-q07}
\end{figure}

\begin{figure}[ht]
	\marginfignote{3}{This plot shows a much broader plane of the final period of a grid of models
		with \(q=0.7,\;\varepsilon=0.1\) and \(f_\text{delta}=0.1\).}
	\centering
	\incplot{w35-mass-5-q07}
\end{figure}

\begin{figure}[ht]
	\marginfignote{3}{This plot shows a much broader plane of the final period of a grid of models
		with \(q=0.7,\;\varepsilon=0.1\) and \(f_\text{delta}=0.3\).}
	\centering
	\incplot{w35-mass-4-q07}
\end{figure}
\end{document}
